
%\section{Einführung}
%\begin{handwritenotes}[Write your thoughts or examples below:]
\section{Definitionen der Stadt}
\begin{enumerate}
    \item \textbf{Statistische Definition}
          \begin{figure}[H]
              \centering
              \includegraphics[width=0.8\linewidth]{"Geografie/Stadtgeografie/Figures/surseePopulation"}
              \caption{Bevölkerungswachstum der Stadt Sursee über die Zeit}
          \end{figure}
          \begin{multicols}{2}
              \notiz{4}
              \columnbreak
              \begin{table}[H]
                  \centering
                  \begin{tabular}{c|c}
                      \textbf{Stadtdefinition} & \textbf{Schwelle} \\\midrule\midrule
                      Island, Dänemark         & 200               \\\midrule
                      Deutschland, Frankreich  & 2'000             \\\midrule
                      Schweiz                  & 10'000            \\\midrule
                      Japan                    & 50'000            \\
                  \end{tabular}
              \end{table}
          \end{multicols}
    \item \textbf{Subjektive Definition}
          \notiz{7}
          \newpage
    \item \textbf{Historische Definition}
          \begin{figure}[H]
              \centering
              \includegraphics[width=.5\linewidth]{"Geografie/Stadtgeografie/Figures/sempach"}
              \caption{Kleinstadt Sempach LU, $\sim$4200 Einwohner}
          \end{figure}
          \begin{figure}[H]
              \centering
              \includegraphics[width=\linewidth]{"Geografie/Stadtgeografie/Figures/histBegriffTemplate"}
              \caption{Dimensionen historischer Städte nach \textcite{hep}}
          \end{figure}
          \begin{multicols}{2}
              \notiz{4}
              \columnbreak
              \begin{table}[H]
                  \textbf{Stadt-Rechte}
                  \begin{enumerate}
                      \item Münzrecht
                      \item Marktrecht
                      \item Befestigungsrecht
                      \item etc.
                  \end{enumerate}
              \end{table}
          \end{multicols}
          \newpage
    \item \textbf{Funktionale Definition}
          \begin{figure}[H]
              \centering
              {\transparent{0.4}\includegraphics[width=0.8\linewidth]{"Geografie/Stadtgeografie/Figures/funktionsraum"}}
              \caption{Funktionale Aspekte einer Stadt (Paris: Aussicht vom \textit{Arc de Triomphe})}
          \end{figure}
          \begin{multicols}{2}
              \notiz{4}
              \columnbreak
              \begin{table}[H]
                  \centering
                  \textbf{7 Grunddaseinsfunktionen}
                  \begin{enumerate}
                      \item Arbeiten
                      \item Wohnen
                      \item Sich erholen
                      \item Sich bilden
                      \item (Sich) versorgen (und entsorgen)
                      \item An gemeinschaftlichen Leben teilnehmen
                      \item Sich fortbewegen
                  \end{enumerate}
              \end{table}
          \end{multicols}
    \item \textbf{Geografische Definition} \begin{multicols}{2}
              \notiz{10}
              \columnbreak

              \begin{figure}[H]
                  \centering
                  \includegraphics[width=\linewidth]{"Geografie/Stadtgeografie/Figures/geoStadtbegriff"}
                  \caption{Geographischer Stadtbegriff nach \textcite{westermann}}
              \end{figure}

          \end{multicols}
\end{enumerate}
\newpage
\section{Städtische Siedlungsformen}
\begin{itemize}
    \item Agglomeration
          \begin{figure}[H]
              \centering
              \includegraphics[width=0.7\linewidth]{"Geografie/Stadtgeografie/Figures/agglomeration"}
              \caption{Agglomerationsgebiete der Schweiz (Ausschnitt), \href{https://www.atlas.bfs.admin.ch/maps/13/map/mapIdOnly/27623_de.html}{Quelle}}
          \end{figure}
          \notiz{4}
    \item Raum mit städtischem Charakter \notiz{4}
    \item Megastadt \notiz{4}
    \item Metropole \notiz{4}
    \item Metropol-Region
          \begin{multicols}{2}
              \notiz{6}
              \columnbreak
              \begin{figure}[H]
                  \centering
                  \includegraphics[width=\columnwidth]{"Geografie/Stadtgeografie/Figures/gza"}
                  \caption{Mitgliedskantone der \enquote{Greater Zurich Area}}
              \end{figure}
          \end{multicols}
    \item Global City \notiz{4}
    \item Megalopolis
          \begin{multicols}{2}
              \notiz{6}
              \columnbreak
              \begin{figure}[H]
                  \centering
                  \includegraphics[width=\columnwidth]{"Geografie/Stadtgeografie/Figures/megalopolis"}
                  \caption{Megalopolis in den USA}
              \end{figure}
          \end{multicols}
\end{itemize}


\begin{myexercise}\label{ex:stadtDef}
    Suchen Sie sich in 3er- oder 4er-Gruppen eine Stadt aus, die Sie im Rahmen der nächsten 7 Doppellektionen vertiefen werden:
    \begin{itemize}
        \item Basel (\& Laufen)
        \item Bern (\& Huttwil)
        \item Genf (\& Allaman)
        \item Andere Stadt (\& Dorf) nach Wunsch
    \end{itemize}
    Erarbeiten Sie ein Plakat auf nächstes Mal, in welcher Sie die unterschiedlichen Definitionen dieser Städte ausarbeiten (die \enquote{Dörfer} in den Klammern werden wir erst später anschauen).

    Stellen Sie Folgendes dar:
    \begin{itemize}
        \item \textbf{Statistische Definition}: Einwohnerzahl, Einwohnerdichte, sowie andere relevante Kennzahlen
        \item \textbf{Historische Definition}: Historische Entstehung, Spezialrechte etc.
        \item \textbf{Subjektive Definition}: Welche städtischen Aspekte verbinden \textit{Sie} mit dieser Stadt? Weshalb finden Sie (nicht), dass es eine Stadt ist?
        \item \textbf{Funktionale Definition}: Suchen Sie ein Luftbild Ihrer Stadt und annotieren Sie darauf die \textbf{Grunddaseinsfunktionen}!
    \end{itemize}

    \textbf{Formales}: Mindestens 2 Bilder, mindestens 7 Grunddaseinsfunktionen.

    \textbf{Nützliche Ressourcen:}
    \begin{itemize}
        \item \href{https://www.atlas.bfs.admin.ch/maps/13/map/mapIdOnly/27621_de.html}{Statistiken des BFS}, sowie weitere Karten zu Bevölkerung (Anzahl, Dichte etc.) auf derselben Webseite.
        \item \href{https://www.lds.geo.uzh.ch/de/Alle-Orte.html}{Städte-Luftbilder des Universität Zürich}
        \item \href{https://hls-dhs-dss.ch/de/articles/000209/2016-11-10/}{Historisches Lexikon der Schweiz, z.B. zur Stadt Bern}
    \end{itemize}

\end{myexercise}
\newpage
\begin{mychallenge}[Ausschnitt aus Buch \enquote{Klett Fundamente} zu Stadtbegriffen]
    \foreach \n in {1,...,3}{
            \begin{figure}[H]
                \centering
                \includegraphics[width=.8\linewidth, page=\n]{"Geografie/Stadtgeografie/Appendix/Der geographische Stadtbegriff.pdf"}
            \end{figure}
        }
\end{mychallenge}

\clearpage
\section{Lernziele}
\begin{todolist}
    \item Ich kann Städte nach statistischen, historischen, funktionalen, geographischen und subjektiven Charakteristiken beschreiben und definieren.
    \item Ich verstehe, dass Städte funktionale Einflussgebiete haben, die weit über das Stadtzentrum hinausgehen und kann beschreiben, inwiefern Städte Grunddaseinsfunktionen für das Umland erfüllen.
    \item Ich kann die Kriterien von Agglomerationen auflisten und verwenden, um städtische Gebiete in eine Kernstadt und eine Agglomeration aufzuteilen.
    \item Ich kann eine Stadt einordnen als Megastadt, Metropole, Megalopolis oder Global City.
\end{todolist}